
\section{Introduction}
\label{sec:intro}

Joint estimation of camera pose, 3D geometry, and 3D motion -- a task known as 4D scene reconstruction 
-- from casually captured images is a fundamental challenge in computer vision with applications in robotics~\citep{xu2024flow}, autonomous driving~\citep{geiger2012we}, and 3D/4D generative AI~\citep{yu20244real, liang2024diffusion4d, xie2024sv4d}. 
Recovering dense 4D information from limited 2D inputs, however, is an inherently ill-posed problem.
While data-driven approach can be a solution, a challenge continues with a scarcity of dense, large-scale 4D training data. 
Synthetic datasets~\citep{Zheng:2023:POA} provide dense, pixel-perfect supervision, yet they often suffer from a significant domain gap and a lack of diversity.
Real-world data~\citep{Jin:2025:S4D} is generally constrained to sparse and noisy annotations, limiting its effectiveness for training robust models. 

Consequently, dynamic scene reconstruction has traditionally centered on slow, test-time optimization pipelines
This approach relies on intermediate 2D signals such as depth and optical flow, which imposes a high computational cost and caps performance at the quality of these inputs.
A recent shift towards feedforward models, including DUST3R~\citep{Wang:2024:DUS}, MonST3R~\citep{Zhang:2025:MST}, and DynaDUSt3R~\citep{Jin:2025:S4D}, has yielded impressive results on individual perception tasks; however, a single, unified feedforward architecture capable of holistically addressing the full spectrum of 2D and 3D perception tasks does not yet exist.
This lack of a unified representation prevents models from exploiting the tightly coupled nature of geometry and motion.

To address these limitations, we introduce \ourmethod{}, a unified feedforward model for dense 4D reconstruction from just two unposed images.
\ourmethod{} predicts an explicit and dense spatio-temporal representation, Dynamic 3D Gaussian Splatting (D-3DGS), in a single forward pass.
This unified representation enables the consistent, integrated estimation of multiple 2D and 3D reconstruction tasks, including depth, scene flow, and optical flow, all from the same underlying model.

Furthermore, our unified representation unlocks two key advantages during training.
First, differentiable rendering provides a dense, self-supervised photometric loss that is independent of sparse or noisy ground-truth labels.
Second, because 3D geometry and motion originate from the same set of Gaussian primitives, their respective losses are tightly coupled in 3D.
This coupling creates a regularizing effect where each supervisory signal helps compensate for imperfections in the other.
As a result, \ourmethod{} achieves state-of-the-art performance on geometry and motion benchmarks, achieving more than $3\times$ lower EPE on Stereo4D and KITTI than competing methods. 
Furthermore, our explicit 4D representation enables novel applications, including high-fidelity spatio-temporal interpolation of appearance, geometry, and motion from a single prediction.

Our primary contributions are:
\begin{itemize}
  \item A unified model for unposed feedforward 4D reconstruction from two images using a dynamic 3D Gaussian Splatting representation.
  \item A robust semi-supervision framework that leverages differentiably rendered outputs to mitigate the scarcity of annotated data.
  \item 4D spatio-temporal interpolation of image, depth, and motion as a new application of the feedforward output.
  \item State-of-the-art performance on 3D geometry and 3D motion benchmark datasets. 
\end{itemize}
